\section{Data and Label Taxonomy}
\label{sec:data}

\subsection{Label taxonomy}

We adopt the ATBench \citep{atbench} three-tier taxonomy directly as our label
space, because it is the most fine-grained publicly available labeled-trajectory
schema and matching it makes joint evaluation trivial. A trajectory's anomaly
label carries:
\begin{itemize}[leftmargin=1.4em,itemsep=0.1em]
  \item \textbf{Risk source} (where the risk originates): user prompt injection,
        tool-description injection, indirect prompt injection, corrupted tool
        feedback, agent hallucination, over-agency, misconfigured permission,
        ambiguous instruction.
  \item \textbf{Failure mode} (what the agent does wrong): unauthorized
        disclosure, data exfiltration, unauthorized action, destructive action,
        privilege escalation, misleading information, instruction for harm,
        resource exhaustion, failed refusal, hallucinated tool use, chain
        hijack, side-effect leak, poisoned memory.
  \item \textbf{Harm category}: privacy, security, financial, physical health,
        fairness, legal, psychological, reputational, system integrity.
\end{itemize}
The label also records \texttt{is\_anomaly}, the \texttt{anomaly\_step} (the
decision step that should be blocked), and a gold \texttt{reason} string used to
supervise the \texttt{<reason>} field and to score reason quality at evaluation.

\subsection{Cross-backend canonicalization}
\label{sec:data:canon}

A central design goal is that one model serves many agent frameworks. We
generate traces from multiple backends --- Claude Code (a subprocess emitting
stream-JSON events) and a custom ReAct loop on the OpenAI Chat Completions API,
with LangGraph as an additional target --- each of which produces a different
raw event format. We fold all of them into the single canonical
\texttt{Trajectory} schema of Section~\ref{sec:method:schema}. Our finding is
that violation patterns defined over the canonical schema transfer across
backends: a monitor trained on one framework's canonicalized traces generalizes
to another's, provided the attack pattern is expressed at the canonical level
rather than at the raw-event level. This is the abstraction that lets TraceGuard
ship as one model usable across frameworks; prior trajectory-anomaly work is
typically single-framework. % TODO: report quantitative cross-backend transfer (train-on-A / test-on-B F1 drop); not yet measured in v0.19.

\subsection{Trace synthesis}

Adversarial traces are generated by running attack scenarios against
weakly-aligned backends (where the attacks succeed) and capturing the resulting
trajectories, complemented by ATBench's labeled set. Attack scenarios span
hand-authored families (direct exfiltration, SSRF, shell injection,
destructive commands, persistence, output manipulation, indirect injection at
several sophistication tiers) and LLM-synthesized extensions of each family; we
find that an LLM-in-the-loop scenario designer scales red-team coverage well
beyond manual authoring (Section~\ref{sec:experiments}, F4). Benign traces are
drawn from multiple sources, including ATBench's benign split and normal agent
runs, so the model sees a realistic mix of safe behavior to learn ``do not stop
here.''

\paragraph{Training set.}
The v0.19 training set comprises $\approx 2{,}157$ labeled trajectories. We
train with full-parameter fine-tuning for 4 epochs at learning rate
$2\times10^{-5}$, \texttt{max\_seq\_len}$=4096$ with left-truncation
(Section~\ref{sec:method:lefttrunc}), and $W_{\textsc{STOP}}=8$ to counter the
$\approx 7{:}1$ \textsc{OK}{:}\textsc{STOP} imbalance. Train and evaluation
splits avoid scenario leakage. % TODO: report exact per-source counts and the train/val/test split sizes.
